% ============================================================================
% EVIDENCE BASE FOR EPISTEMIC PROMPTING
% ============================================================================
% Status: FINAL - Scientific legitimation of ESL methodology
% Date: 2025-12-28
% Purpose: Ground ESL in established cognitive science literature
% ============================================================================

\documentclass[11pt,a4paper]{article}
\usepackage[utf8]{inputenc}
\usepackage[T1]{fontenc}
\usepackage{amsmath,amssymb}
\usepackage{booktabs}
\usepackage{longtable}
\usepackage{tcolorbox}
\usepackage[margin=1in]{geometry}
\usepackage{natbib}

\newtcolorbox{keypoint}{
  colback=green!5!white,
  colframe=green!50!black,
  sharp corners
}

\title{\textbf{Evidence Base for Epistemic Prompting}\\
\Large Scientific Foundations of the ESL Framework}
\author{FehrAdvice \& Partners AG\\BEATRIX Research Group}
\date{December 2025}

\begin{document}
\maketitle

% ============================================================================
\section{Executive Summary}
% ============================================================================

\begin{keypoint}
The prompting strategy used in the Activation Challenge operationalizes established findings from cognitive science, metacognition, decision research, and Human--AI interaction.

\textbf{It is not a new assumption, but a formal bundling of known effects.}
\end{keypoint}


% ============================================================================
\section{Evidence Domain 1: Knowledge Possession vs. Application}
% ============================================================================

\subsection{Finding}

Humans (and systems) can:
\begin{itemize}
    \item Possess knowledge
    \item Without using it
\end{itemize}

\subsection{Key Evidence}

\begin{itemize}
    \item \textbf{Availability vs. Accessibility} \citep{tulving1966availability}
    \item \textbf{Inert Knowledge Problem} \citep{whitehead1929aims}
    \item \textbf{Transfer Failure} in learning and expertise research
\end{itemize}

\subsection{Relevance for Prompting}

$\rightarrow$ Explicit retrieval conditions increase \emph{usage}, not possession.

\begin{keypoint}
ESL prompt = formal externalization of this separation.
\end{keypoint}


% ============================================================================
\section{Evidence Domain 2: Metacognitive Calibration \& Overconfidence}
% ============================================================================

\subsection{Finding}

Humans overestimate knowledge when:
\begin{itemize}
    \item No explicit verification rules are present
    \item Plausibility is confused with evidence
\end{itemize}

\subsection{Key Evidence}

\begin{itemize}
    \item \textbf{Overconfidence Bias} \citep{fischhoff1977knowing}
    \item \textbf{Illusion of Explanatory Depth} \citep{rozenblit2002misunderstood}
    \item \textbf{Metacognitive Monitoring Failures}
\end{itemize}

\subsection{Relevance for Prompting}

$\rightarrow$ Explicit ``stop rules'' reduce false-positive judgments.

\begin{keypoint}
ESL rules (``similar $\neq$ identical'') are exactly such stop rules.
\end{keypoint}


% ============================================================================
\section{Evidence Domain 3: Dual-Process Models}
% ============================================================================

\subsection{Finding}

\begin{itemize}
    \item Intuitive processes dominate by default
    \item Analytical processes must be explicitly activated
\end{itemize}

\subsection{Key Evidence}

\begin{itemize}
    \item \textbf{Thinking, Fast and Slow} \citep{kahneman2011thinking}
    \item \textbf{Dual-process theories} \citep{evans2013dual}
    \item \textbf{Cognitive Reflection Test (CRT)} \citep{frederick2005cognitive}
\end{itemize}

\subsection{Relevance for Prompting}

$\rightarrow$ Without coercion, the system remains in fluency/plausibility mode.

\begin{keypoint}
ESL prompt forces System-2-like verification.
\end{keypoint}


% ============================================================================
\section{Evidence Domain 4: Decision Rules \& Explicit Criteria}
% ============================================================================

\subsection{Finding}

Decisions become more robust when:
\begin{itemize}
    \item Criteria are made explicit
    \item Implicit heuristics are forbidden
\end{itemize}

\subsection{Key Evidence}

\begin{itemize}
    \item \textbf{Fast-and-frugal heuristics} \citep{gigerenzer1999simple}
    \item \textbf{Decision hygiene} \citep{kahneman2021noise}
    \item \textbf{Checklist effects} in medicine and aviation
\end{itemize}

\subsection{Relevance for Prompting}

$\rightarrow$ Explicit criteria reduce variance and bias.

\begin{keypoint}
Baseline prompt = checklist without coercion\\
ESL prompt = checklist with hard exclusion rules
\end{keypoint}


% ============================================================================
\section{Evidence Domain 5: Instruction Following $\neq$ Knowledge Use}
% ============================================================================

\subsection{Finding (LLM-Specific)}

LLMs:
\begin{itemize}
    \item Follow instructions syntactically
    \item Without epistemic guarantee
\end{itemize}

\subsection{Key Evidence}

\begin{itemize}
    \item Instruction-following $\neq$ factual reliability
    \item \textbf{Faithfulness vs. Fluency} \citep{ji2023survey}
    \item Hallucination under instruction compliance
\end{itemize}

\subsection{Relevance for Prompting}

$\rightarrow$ Instructions must be \emph{epistemic}, not just procedural.

\begin{keypoint}
ESL is \textbf{epistemic prompting}, not task prompting.
\end{keypoint}


% ============================================================================
\section{Evidence Domain 6: Label Precision}
% ============================================================================

\subsection{Finding}

\begin{itemize}
    \item Label accuracy is central for knowledge access
    \item Semantic proximity creates false confidence
\end{itemize}

\subsection{Key Evidence}

\begin{itemize}
    \item \textbf{Label leakage} in machine learning
    \item \textbf{Concept--label dissociation}
    \item \textbf{Terminological precision} in science
\end{itemize}

\subsection{Relevance for Prompting}

$\rightarrow$ ``Moral Licensing'' $\neq$ ``Moral Licensing Spillover Effect''

\begin{keypoint}
ESL makes this separation explicit and verifiable.
\end{keypoint}


% ============================================================================
\section{Mapping to $\alpha$-Theory}
% ============================================================================

\begin{table}[h]
\centering
\caption{Evidence domains mapped to theoretical constructs}
\begin{tabular}{ll}
\toprule
\textbf{Evidence Domain} & \textbf{Correspondence in $\alpha$-Theory} \\
\midrule
Inert knowledge & $M_{\text{latent}} \neq A$ (knowledge $\neq$ application) \\
Metacognitive stop rules & Activation of $\alpha$ \\
Dual-process & Fluency mode vs. Activation mode \\
Decision rules & ESL calibration rules \\
Instruction failures & Prompt $\neq$ Knowledge \\
Label precision & $K \leq M$ principle \\
\bottomrule
\end{tabular}
\end{table}

\begin{keypoint}
Our prompting is a formal operationalization of all these findings.
\end{keypoint}


% ============================================================================
\section{Paper-Ready Paragraph}
% ============================================================================

\begin{quote}
Our prompting strategy is grounded in established findings from cognitive science and decision research, including the distinction between knowledge availability and accessibility \citep{tulving1966availability}, failures of metacognitive calibration \citep{rozenblit2002misunderstood}, and dual-process models of reasoning \citep{kahneman2011thinking}. Explicit epistemic rules and exclusion criteria have been shown to reduce overconfidence and heuristic-driven errors in human decision-making \citep{kahneman2021noise}. We translate these principles into a formal prompting framework that does not add information, but constrains the conditions under which latent knowledge may be applied.
\end{quote}


% ============================================================================
\section{What We Do NOT Claim}
% ============================================================================

\begin{itemize}
    \item[$\times$] That prompting creates knowledge
    \item[$\times$] That ESL ``trains'' the model
    \item[$\times$] That rules always help
\end{itemize}

\begin{itemize}
    \item[$\checkmark$] \textbf{Only:} Prompting can formalize known cognitive control mechanisms---or fail.
\end{itemize}


% ============================================================================
\section{Summary}
% ============================================================================

\begin{keypoint}
\textbf{ESL prompting is evidence-based because it formalizes established cognitive mechanisms for controlling knowledge application and transfers them to LLMs.}
\end{keypoint}


% ============================================================================
% REFERENCES
% ============================================================================

\bibliographystyle{apalike}
\begin{thebibliography}{99}

\bibitem{tulving1966availability}
Tulving, E., \& Pearlstone, Z. (1966). Availability versus accessibility of information in memory for words. \emph{Journal of Verbal Learning and Verbal Behavior}, 5(4), 381--391.

\bibitem{whitehead1929aims}
Whitehead, A. N. (1929). \emph{The Aims of Education and Other Essays}. Macmillan.

\bibitem{fischhoff1977knowing}
Fischhoff, B., Slovic, P., \& Lichtenstein, S. (1977). Knowing with certainty: The appropriateness of extreme confidence. \emph{Journal of Experimental Psychology: Human Perception and Performance}, 3(4), 552--564.

\bibitem{rozenblit2002misunderstood}
Rozenblit, L., \& Keil, F. (2002). The misunderstood limits of folk science: An illusion of explanatory depth. \emph{Cognitive Science}, 26(5), 521--562.

\bibitem{kahneman2011thinking}
Kahneman, D. (2011). \emph{Thinking, Fast and Slow}. Farrar, Straus and Giroux.

\bibitem{evans2013dual}
Evans, J. St. B. T., \& Stanovich, K. E. (2013). Dual-process theories of higher cognition: Advancing the debate. \emph{Perspectives on Psychological Science}, 8(3), 223--241.

\bibitem{frederick2005cognitive}
Frederick, S. (2005). Cognitive reflection and decision making. \emph{Journal of Economic Perspectives}, 19(4), 25--42.

\bibitem{gigerenzer1999simple}
Gigerenzer, G., Todd, P. M., \& the ABC Research Group. (1999). \emph{Simple Heuristics That Make Us Smart}. Oxford University Press.

\bibitem{kahneman2021noise}
Kahneman, D., Sibony, O., \& Sunstein, C. R. (2021). \emph{Noise: A Flaw in Human Judgment}. Little, Brown Spark.

\bibitem{ji2023survey}
Ji, Z., et al. (2023). Survey of hallucination in natural language generation. \emph{ACM Computing Surveys}, 55(12), 1--38.

\end{thebibliography}

\end{document}
